\section{Conclusions}

Empirical valence bond calculations show that the catalytic advantage of selenocysteine in GPX6 is not a fixed local property of the reactive residue, but is modulated by the surrounding protein scaffold. Introducing selenocysteine into mouse GPX6 lowers the proton-transfer barrier more strongly than the reciprocal cysteine substitution raises it in human GPX6, demonstrating that the same Sec/Cys exchange has different energetic consequences in different sequence backgrounds.

The mutational pathway analysis further shows that the transition between human and mouse GPX6 is governed by epistatic constraints. The evaluated paths are diverse, and the triangular sequence projection does not support a simple picture of strict hysteresis. However, the EVB landscapes reveal a recurring energetic bias in how substitutions convert one catalytic background into the other. Some substitutions appear compatible with both Sec- and Cys-containing states, whereas others are more strongly contingent on the catalytic residue, providing a way to narrow plausible evolutionary orders. GPX6 therefore illustrates how a chemically meaningful residue replacement can become embedded within lineage-specific compensatory networks that reshape transition-state stabilization, even when the derived Cys state may have been favored by selection for a distinct function. This framework should be useful for studying other Sec/Cys transitions and for designing enzyme-engineering strategies that account for catalytic chemistry and sequence background together.
